\section{Experimental Protocol}
\label{sec:experiments}

\paragraph{Objectives.}
Using the matched comparison of Section~\ref{sec:study}, our analysis
establishes three things left unresolved by aggregate efficient-decoder
evaluation: how large and how bidirectional the full decoder's benefit is, where
it is concentrated anatomically, and whether the beneficial regions can be
identified at test time from the efficient model's uncertainty.
Section~\ref{sec:results} reports these in turn, preceded by the
controlled-reimplementation baselines.

\paragraph{Dataset and split.}
The primary task is the 13-organ BTCV/Synapse abdominal CT benchmark with 18
training and 12 validation volumes.  The evaluated structures are spleen, right
and left kidney, gallbladder, esophagus, liver, stomach, aorta, inferior vena
cava, portal and splenic veins, pancreas, and right and left adrenal glands.
Images are intensity-windowed and resampled following the public \method{}
pipeline; training uses $96^3$ patches.  The same validation subjects are paired
across every model comparison.

\paragraph{Models and training.}
\fullmodel{} is the public full 3D UX-Net~\cite{lee20233duxnet} configuration
with feature widths $[48,96,192,384]$; \effimodel{} replaces its decoder with the
\method{} decoder (48 channels, additive skip aggregation, resolution factor
two).  Both are trained for 45k iterations with a Dice--cross-entropy loss and
AdamW.  We evaluate one \fullmodel{} run and two independently initialized
\effimodel{} runs.  Sliding-window inference uses a $96^3$ window, batch size
four, overlap 0.7, and constant blending.  We report Dice and 95th-percentile
Hausdorff distance (HD95), together with parameters, multiply--accumulate
operations (MACs), latency, and peak inference memory.

\paragraph{Reimplementation status.}
Our data were reconstructed from the public TransUNet/Synapse release, whose
NIfTI files carry identity affine metadata, whereas \method{} refers to a
processed \texttt{btcv\_trns} dataset that is not distributed.  We therefore
present the experiment as a \emph{controlled reimplementation}, not an exact
reproduction, and disclose the discrepancy from published numbers in full.
Because every observation is relational---matched full versus efficient on
identical local data---a residual reproduction gap for any one cell affects at
most the estimated \emph{magnitude} of that cell's opportunity, which we bound
with net (not positive-only) transitions and subject-level intervals; it does
not change the direction of the finding.  We verified that this gap is not a
precision artifact: re-evaluating the same weights in FP32 and BF16 gives
identical Dice.

\paragraph{Architecture generality.}
We organize architecture coverage as a predeclared ladder over the backbones for
which \method{} actually defines a decoder pair, rather than treating one extra
model as proof of universality.  The primary causal comparison is the matched
3D UX-Net pair (large-kernel CNN); a matched Swin UNETR~\cite{tang2022swinunetr}
pair tests transfer to a hierarchical Transformer, with
SwinUNETR-V2~\cite{he2023swinunetrv2} and MedNeXt~\cite{roy2023mednext}
available as optional deeper pairs.  We deliberately exclude
UNETR~\cite{hatamizadeh2022unetr} and nnU-Net~\cite{isensee2021nnunet}: \method{}
builds an efficient decoder only where a wide, high-resolution decoder
bottleneck exists, and pairs one with neither (UNETR is already lightweight;
nnU-Net self-configures its decoder).  A matched transition comparison is
therefore undefined for them, and an unofficial \method{} variant would be an
ungrounded extension rather than a control.

\paragraph{Dataset generality.}
The original \method{} evaluation spans FeTA~\cite{payette2021feta}, BTCV, and
the ten Medical Segmentation Decathlon (MSD) tasks~\cite{antonelli2022msd}.  We
preselect three to four representative tasks: BTCV abdominal CT as the primary
multi-organ task; FeTA fetal-brain MRI for modality and anatomy shift; MSD
Task01 BrainTumour for multimodal lesion segmentation; and either MSD Task06
Lung or Task08 HepaticVessel as a sparse-lesion or thin-structure stress test.
Every dataset uses the same positive, negative, and net transition definitions,
concentration analysis, opportunity curve, and subject-level intervals, and the
set is frozen before inspecting results.

\todo{Insert the matched Swin UNETR replication (opt.\ SwinUNETR-V2/MedNeXt) and
the predeclared FeTA/MSD results; all currently completed numbers are 3D UX-Net
on BTCV.}
